\subsection{The remaining large minimal bands}
\label{A:large-band-target-proof}

The argument in this subsection is specific to the projective-plane
target. It avoids the stronger inherited assertion about arbitrary
modules over four-dimensional contracted monoid algebras. Throughout,
algebras and modules in this subsection are finite dimensional over
$\mathbb Q$, and all modules are unital.

\begin{lemma}[Proved: local-algebra facts]
\label{A:local-target-facts}
Let $A$ be local, with radical $J$ and residue division algebra $D$.
For nonzero right and left modules $W,V$, respectively, there is a
surjection
\[
 W\otimes_A V\twoheadrightarrow
 (W/WJ)\otimes_D(V/JV).
\]
Consequently, if $\dim_{\mathbb Q}(W\otimes_AV)=1$, then
$D=\mathbb Q$ and both $W,V$ are cyclic. Moreover, for a
nonprojective finite right module $W$,
$\operatorname{Tor}_i^A(W,D)\ne0$ for every $i\geq1$;
the analogous assertion holds with the sides interchanged.
\end{lemma}
\begin{proof}
Quotienting both tensor factors gives the displayed surjection.
If the tops have dimensions $a,b$ over $D$, its target has
rational dimension $ab\dim_{\mathbb Q}D$. Nonzero finite
modules have nonzero tops by nilpotence of $J$. A one-dimensional
tensor product therefore forces $a=b=1$ and $D=\mathbb Q$.
A lift of a top generator generates the whole module, again
by nilpotence of $J$.

Finite projectives over a local finite algebra are free: choose
a basis of their tops, lift it to a surjection from a free
module, and split the surjection; the complementary projective
has zero top and hence is zero. Construct a minimal free
resolution by taking top bases at each stage. Its differentials
have entries in $J$. Such a resolution cannot end with a
nonzero free module in positive degree. Indeed a nonzero
element in the last nonzero radical power is annihilated by
$J$ on both sides and supplies a nonzero vector in the kernel
of that last differential. Tensoring the nonterminating minimal
resolution with $D$ makes every differential zero, proving
nonvanishing in every positive degree.
\end{proof}

\begin{lemma}[Proved: nonlocal algebras through dimension four]
\label{A:nonlocal-four}
If $A$ is nonlocal and $\dim_{\mathbb Q}A\leq4$, then
$\operatorname{Tor}_2^A(W,V)\ne0$ implies nonzero Tor in
arbitrarily high degrees.
\end{lemma}
\begin{proof}
Write $S=A/J$. Semisimple algebras have zero positive Tor.
We record the elementary small-dimension analysis, including
the nonsplit possibilities. A nonzero radical bimodule
incident with a field factor of degree $d$ has dimension at
least $d$, since it is a unital vector space over that field.
A matrix factor of dimension four similarly cannot act on
a one-dimensional radical. Thus a nonsplit nonlocal algebra
of dimension at most four with nonzero radical can only have
a quadratic field factor, a rational factor and a
one-dimensional radical supported on the rational factor.
It is the product of that field and the dual numbers.
This has the asserted Tor behavior.

In the remaining nonsemisimple cases $S=\mathbb Q^r$,
$r\geq2$. Lift the primitive orthogonal idempotents. Such
lifts can be obtained successively modulo the radical powers;
for a square-zero error $\delta=p^2-p$, the correction
$p+(1-2p)\delta$ is idempotent. Here $\dim J\leq2$.
If $J^2\ne0$, then $r=2$ and $J/J^2,J^2$ are
one dimensional. A generator of $J/J^2$ must be a loop:
a single off-diagonal component cannot multiply with itself.
Lifting that loop gives $x^2\ne0$, $x^3=0$, and the algebra
is $\mathbb Q[x]/(x^3)\times\mathbb Q$.

For $\mathbb Q[x]/(x^m)$, $m\geq2$, Jordan decomposition
expresses every finite module as a sum of $A/(x^a)$.
For $a<m$ its free resolution alternates multiplication by
$x^a$ and $x^{m-a}$. Tensoring with $A/(x^b)$, $b<m$,
has positive-degree homology dimension
\[
 \min(a,b)+\min(m-a,b)-b
 =\min\{a,b,m-a,m-b\}>0.
\]
Free summands contribute no positive Tor. This proves the
claim for the truncated-polynomial products above.

It remains that $J^2=0$. The idempotent decomposition makes
$A$ the radical-square-zero algebra of a directed quiver,
with $r$ vertices and $d=\dim J$ arrows, where $r+d\leq4$.
First syzygies $\Omega W,\Omega V$ are semisimple because
they lie in radicals of projective covers. Dimension shifting
on both sides gives
\[
 \operatorname{Tor}_n^A(W,V)
   \cong\operatorname{Tor}_{n-2}^A(\Omega W,\Omega V)
 \quad(n\geq2).
\]
For $n=2$ this follows from the long exact sequence and the
zero map $\Omega W\otimes_A\Omega V\to
\Omega W\otimes_A P_V$: the second syzygy factor maps into
$JP_V$, while $\Omega WJ=0$.
Nonzero degree-two Tor therefore means that a vertex simple
occurs in both syzygies. That vertex has an incoming and an
outgoing arrow. With at most two arrows and $r+d\leq4$,
it lies on a loop or a two-cycle; a line through three
distinct vertices would require at least five dimensions.
The minimal resolution of a vertex simple in a
radical-square-zero algebra has basis the directed paths
starting at that vertex. Repeating this cycle gives positive
Tor of the two matching simples in unbounded degrees.
Direct sums and the displayed shift finish the proof.
\end{proof}

\begin{lemma}[Proved: a small local summand cannot supply the target]
\label{A:local-four-small-target}
Let $A$ be local of dimension four. There are no nonzero modules
$\dim V\leq2$, $\dim W\leq5$ satisfying
\[
 \bigl(\dim\operatorname{Tor}_0,
       \dim\operatorname{Tor}_1,
       \dim\operatorname{Tor}_2,
       \dim\operatorname{Tor}_3\bigr)=(1,0,1,0).
\]
\end{lemma}
\begin{proof}
Lemma~\ref{A:local-target-facts} makes the residue field
$\mathbb Q$ and both modules cyclic. Hence $\dim W\leq4$.
If either module has dimension one, its positive Tor with a
nonprojective opposite module never vanishes. If $\dim W=4$,
the cyclic cover $A\to W$ is an isomorphism, so positive Tor
vanishes. If $\dim W=3$, its first syzygy is the
one-dimensional simple module, and
$\operatorname{Tor}_3(W,V)=\operatorname{Tor}_2(\mathbb Q,V)
\ne0$ because $V$ is nonprojective. Only
$\dim W=\dim V=2$ remains.

Put $I=\Omega V$. The cyclic cover has kernel of dimension
two. Since $\operatorname{Tor}_1(W,V)=0$,
$W\otimes I$ embeds in $W$, with image the kernel of
$W\to W\otimes V$; its dimension is one. The quotient
onto $\operatorname{top}W\otimes\operatorname{top}I$
therefore makes $I$ cyclic. Put $L=\Omega^2V$, also of
dimension two. If $L$ were not cyclic it would be the sum
of two simples, giving
$\operatorname{Tor}_3(W,V)=\operatorname{Tor}_1(W,L)\ne0$.
Thus $L$ is cyclic. Its syzygy $I'=\Omega^3V$ has dimension
two. The equality $\operatorname{Tor}_1(W,L)=0$ embeds
$W\otimes I'$ into $WJ$, of dimension one, so $I'$ is
cyclic as well.

The first four projectives in the minimal resolution of $V$
are consequently $A$. After tensoring with $W$, each has
dimension two. Since $\operatorname{Tor}_0$ has dimension
one, the first differential has rank one. Vanishing of
$\operatorname{Tor}_1$ makes the second differential rank
one. Dimension one in $\operatorname{Tor}_2$ makes the
third differential zero. But the fourth differential has
image in $WJ$, of dimension one, so it cannot cover the
two-dimensional kernel of the third. This contradicts
$\operatorname{Tor}_3=0$.
\end{proof}

\begin{lemma}[Proved: the nonlocal five-dimensional alternatives needed here]
\label{A:nonlocal-five-target}
Let $A$ be nonlocal of dimension five. Either nonzero
$\operatorname{Tor}_2(W,V)$ persists in arbitrarily high
degrees for all finite module pairs, or $A$ is one of the
following types:
\begin{enumerate}
\item $B\times\mathbb Q$ with $B$ a four-dimensional local algebra;
\item the radical-square-zero algebra of the directed line on
three vertices;
\item the two-vertex algebra with basis $e_1,e_2,a,b,c$,
$a\in e_1Ae_2$, $b\in e_2Ae_1$, $ab=c\ne0$, $ba=0$,
and $cJ=Jc=0$.
\end{enumerate}
\end{lemma}
\begin{proof}
We again use $S=A/J$. Nonsplit semisimple factors require
their degree in every incident radical component. The only
new nonsemisimple possibility relative to
Lemma~\ref{A:nonlocal-four} is a quadratic field and a
rational factor with a two-dimensional radical. A component
joining the two factors is a single triangular arrow and
gives a hereditary algebra: its bimodule is projective on
both sides over semisimple fields, and the standard
one-arrow projective resolution has length one.
A radical component on the quadratic factor gives a local
square-zero algebra over that division residue algebra;
its first syzygies are semisimple and
Lemma~\ref{A:local-target-facts} gives Tor persistence.
A radical supported on the rational factor gives an algebra
of dimension at most three, already covered by the
square-zero or truncated-polynomial arguments. Larger field
factors and matrix factors leave at most one radical
dimension and give no additional case. Separable field
factors split across the nilpotent ideal by lifting a
primitive element and correcting its minimal polynomial;
its derivative is invertible, so the usual successive
Newton correction applies.

Now assume $S=\mathbb Q^r$, $r\geq2$, and lift the
orthogonal idempotents. If $r\geq4$, at most one radical
arrow remains and the preceding argument applies. If
$r=3$ and $J^2\ne0$, a sole loop generates the radical,
giving $\mathbb Q[x]/(x^3)\times\mathbb Q^2$.
If $r=3$ and $J^2=0$, there are two arrows. The
semisimple-syzygy proof of Lemma~\ref{A:nonlocal-four}
applies unless a vertex has an entering and a leaving
arrow but lies on no cycle. The only exception is the
directed line on three vertices, giving alternative (2).

For $r=2$ we have $\dim J=3$. If $J^2=0$, every vertex
with both entering and leaving arrows lies on a loop or
a two-cycle, so the same proof gives persistence.
If $\dim J^2=2$, the single generator of $J/J^2$ is a
loop and gives $\mathbb Q[x]/(x^4)\times\mathbb Q$,
which is covered by (1). Suppose $\dim J^2=1$; then
$J^3=0$ and there are two arrows modulo $J^2$.
Two parallel cross arrows cannot generate a nonzero
square. A loop at each vertex gives a product of
$\mathbb Q[x]/(x^3)$ and $\mathbb Q[y]/(y^2)$.
Two opposite cross arrows give (3), because their
composites lie at distinct vertices and exactly one
can span the one-dimensional square. Two loops at
one vertex give (1).

The remaining configuration, up to opposition, consists
of $x\in e_1Je_1$ and $a\in e_1Je_2$.
The potentially nonzero products $x^2$ and $xa$
lie in different idempotent components, so precisely
one spans $J^2$.
If $x^2=c$ and $xa=0$, write $P_i=e_iA$ and
$T_i$ for right simples. Then
\[
 \operatorname{rad}P_1=L\oplus P_2,\quad
 L=\langle x,c\rangle,\quad P_2=T_2,\qquad
 \Omega L=T_1\oplus P_2,\quad
 \Omega T_1=L\oplus P_2.
\]
Every first right syzygy is a sum of $L,T_1,P_2$:
the cross arrow acts trivially there, and on the
first-vertex component this is Jordan decomposition
for a square-zero operator. Its nonprojective
summands have positive Tor periodic with period two.
If $xa=c$ and $x^2=0$, the same decomposition of
$\operatorname{rad}P_1$ holds, but now $\Omega L=L$.
On $L$ the cross arrow is an isomorphism between
its two one-dimensional vertex spaces. Every submodule
of a sum of $L$ and $P_2$ has injective cross-arrow
map and splits as a sum of those two modules.
Thus first syzygies have positive Tor periodic with
period one. In both cases a nonzero degree-two Tor
group recurs indefinitely. Opposition interchanges
the two tensor factors. These cases exhaust the
two-arrow configurations and prove the alternatives.
\end{proof}

\begin{proposition}[Proved: the $3\times6$ band with five outside elements]
\label{A:three-by-six-CP2}
A finite monoid with a $3\times6$ rectangular minimal band
and exactly five outside elements cannot have classifying
space $\CP^2$. The same holds for its opposite.
\end{proposition}
\begin{proof}
Let $A=\mathbb QM/\mathbb QK$. Its dimension is five;
the row and column augmentation modules have dimensions
$\dim V=2$, $\dim W=5$. The band Tor formula would give
\[
 \dim\operatorname{Tor}_n^A(W,V)=
 \begin{cases}1,&n=0,2,\\0,&n=1\text{ or }n\geq3.\end{cases}
\]
If $A$ is local, Lemma~\ref{A:local-target-facts}
makes $W$ cyclic. Its cover $A\to W$ has equal
source and target dimensions, hence is an isomorphism,
contradicting positive Tor. In the nonlocal case
Lemma~\ref{A:nonlocal-five-target} leaves its three
exceptional alternatives. For $B\times\mathbb Q$,
positive Tor comes entirely from the $B$ summands;
their tensor product is nonzero and hence has dimension
one, since the total tensor dimension is one.
Lemma~\ref{A:local-four-small-target} applies to
those summands and excludes this case.

For the line algebra write $a\in e_1Ae_2$,
$b\in e_2Ae_3$, $ab=0$.
Splitting the image of $b$ inside the kernel of $a$
decomposes every left module into the three simples
$S_1,S_2,S_3$ and the two arrow modules $Ae_2,Ae_3$;
dually the right indecomposables are
$T_1,T_2,T_3,e_1A,e_2A$.
The two nontrivial right resolutions are
\[
 0\to e_3A\xrightarrow{b}e_2A\xrightarrow{a}e_1A
   \to T_1\to0,\qquad
 0\to e_3A\xrightarrow{b}e_2A\to T_2\to0.
\]
Thus degree-two Tor one requires one $T_1$ in $W$
and one $S_3$ in $V$. The remaining dimension of
$V$ must be $S_1$: $S_2$ would give nonzero degree-one
Tor, and a second $S_3$ would give degree-two Tor
dimension at least two. Degree-one vanishing excludes
$T_2$ from $W$. Hence
\[
 V=S_1\oplus S_3,\qquad
 W=T_1\oplus T_3^{\oplus s}\oplus(e_1A)^{\oplus p}
                    \oplus(e_2A)^{\oplus q},\quad
 s+2p+2q=4.
\]
The tensor dimension is $1+s+p$, so it equals one
only for $(s,p,q)=(0,0,2)$.
Every balanced functional for this pair is supported
on the single $T_1\times S_1$ coordinate and has
matrix rank one. Rank is invariant under changes
of module bases; therefore the actual band balanced
form also has rank one. Lemma~\ref{A:rank-one-square}
contradicts the required nonzero square.

For the two-vertex algebra, the left indecomposables
are $S_1,S_2,L=Ae_1/\mathbb Qc,Ae_1,Ae_2$, of
dimensions $1,1,2,3,2$, and the right ones are
$T_1,T_2,R=e_1A/\mathbb Qc,e_1A,e_2A$.
To see completeness, split off the length-three
chains representing a basis of $\operatorname{im}(ab)$.
On the complementary representation both composites
vanish; splitting the two arrow images inside the
opposite kernels leaves length-two chains and simples.
The relevant right resolutions are
\[
\begin{split}
 &0\to e_2A\xrightarrow{a}e_1A\to T_1\to0,\\
 &0\to e_2A\xrightarrow{a}e_1A\xrightarrow{b}e_2A
       \to T_2\to0,\\
 &0\to e_2A\xrightarrow{a}e_1A\xrightarrow{c}e_1A
       \to R\to0.
\end{split}
\]
Evaluating them on the five possible two-dimensional
left modules shows that only $V=L$ can have degree-one
Tor zero and degree-two Tor one: two copies of a simple
give an even degree-two dimension, $S_1\oplus S_2$
forces nonzero degree-one Tor, and $Ae_2$ gives zero
degree-two Tor. With $V=L$, one must have exactly
one $T_2$, and neither $T_1$ nor $R$ can occur in
$W$. Thus
$W=T_2\oplus(e_1A)^{\oplus p}\oplus(e_2A)^{\oplus q}$
with $3p+2q=4$, forcing $p=0,q=2$.
But $T_2\otimes L=0$ and each $e_2A\otimes L$
has dimension one, so the tensor dimension is two,
contradicting the target.

These small calculations are also regenerated from
literal multiplication tables and the displayed
resolutions by
\path{monoid/code/five_dimensional_pairs.py}.
Its complete module-pair output is
\path{monoid/certificates/five_dimensional_pairs.json}.
The proof above establishes their decomposition
scope and uses no list of monoid multiplicative bases.
\end{proof}

\begin{corollary}[Proved: the other large-band shapes through order twenty-three]
\label{A:large-bands-through23}
No finite monoid of order at most twenty-three with
rectangular minimal band of shape $3\times6$, $3\times7$
or $4\times5$ has type $\CP^2$. The same holds after
opposition.
\end{corollary}
\begin{proof}
Let $d=|M\setminus K|=\dim A$. For $3\times6$,
either $d=5$, treated by Proposition~\ref{A:three-by-six-CP2},
or $d\leq4$. For $3\times7$ we have $d\leq2$;
for $4\times5$, $d\leq3$. In each of the latter
cases $A$ has dimension at most four and one of
$V,W$ has dimension strictly greater than $d$.
If $A$ is local, tensor dimension one would make
both modules cyclic by Lemma~\ref{A:local-target-facts},
which is impossible for a module of dimension greater
than its algebra. If $A$ is nonlocal,
Lemma~\ref{A:nonlocal-four} contradicts the target's
vanishing Tor in all degrees above two.
\end{proof}
